\documentclass[12pt,a4paper]{article}

\usepackage[utf8]{inputenc}
\usepackage[T1]{fontenc}
\usepackage[french]{babel}
\usepackage{lmodern}
\usepackage{amsmath,amssymb}
\usepackage{geometry}
\usepackage{microtype}
\usepackage{hyperref}
\usepackage{enumitem}
\geometry{top=2.5cm,bottom=2.5cm,left=2.5cm,right=2.5cm}
\hypersetup{hidelinks,pdftitle={Documentation Hippocampe HCANN}}

\title{Documentation du Bloc Hippocampe}
\author{Projet HCANN-Agent}
\date{\today}

\begin{document}
\maketitle

\section{Bloc Hippocampe (\texttt{hippocampus.py})}

\subsection{Objectif du bloc}
Le bloc hippocampique implémente un \textit{scaffold} de type Vector-HaSH, dédié à l'encodage rapide des épisodes. Il combine :
\begin{itemize}
    \item une intégration de mouvement via des modules de \textit{grid cells} ;
    \item une projection fixe vers un espace hippocampique de grande dimension ;
    \item une dynamique d'attracteur pour stabiliser les états internes.
\end{itemize}

\subsection{Inspiration neuroscientifique}
Le design est inspiré de propriétés fonctionnelles de la formation hippocampique :
\begin{itemize}
    \item \textbf{Grid cells} (cortex entorhinal médian) : codage périodique des transitions latentes ;
    \item \textbf{Séparation structure--contenu} : dissociation entre dynamique interne du support (\textit{scaffold}) et contenu épisodique ;
    \item \textbf{Attracteurs hippocampiques} : stabilisation des traces mnésiques face au bruit.
\end{itemize}
L'approche reste algorithmique : elle vise la plausibilité fonctionnelle et la robustesse computationnelle, sans simulation biophysique fine.

\subsection{1) Module de cellules de grille}
Chaque module de période $P$ maintient une phase 2D $\boldsymbol{\phi}_t$ et applique :
\begin{equation}
\boldsymbol{\phi}_{t+1} = (\boldsymbol{\phi}_t + \mathbf{u}_t)\bmod P,
\end{equation}
où $\mathbf{u}_t\in\mathbb{R}^2$ représente la vitesse (ou transition latente).

Le codage périodique d'un module est défini par :
\begin{equation}
g(\boldsymbol{\phi})=
\left[
\sin\!\left(\tfrac{2\pi}{P}\boldsymbol{\phi}\right),
\cos\!\left(\tfrac{2\pi}{P}\boldsymbol{\phi}\right)
\right]\in\mathbb{R}^{4}.
\end{equation}

\paragraph{Détail d'implémentation}
Le code adapte dynamiquement la taille batch de \texttt{phase} avant l'addition avec $\mathbf{u}_t$, afin d'éviter les incohérences de dimensions entre mini-lots successifs.

\subsection{2) Scaffold hippocampique}
Avec $M$ modules, l'état de grille concaténé est :
\begin{equation}
\mathbf{g} = [g_1,\dots,g_M] \in \mathbb{R}^{4M}.
\end{equation}

La projection hippocampique fixe (\texttt{W\_grid\_to\_hpc}) est ensuite appliquée :
\begin{equation}
\mathbf{h}_0 = \mathrm{ReLU}(W_{g\to h}\mathbf{g}-\theta_h),
\end{equation}
où $W_{g\to h}$ est initialisé aléatoirement puis figé, et $\theta_h=\texttt{hpc\_threshold}$.

\subsection{3) Dynamique d'attracteur}
L'état hippocampique est raffiné par itérations :
\begin{equation}
\mathbf{h}^{(k+1)} = \tanh\!\left(\mathbf{h}^{(k)} W_A\right), \quad k=0,\dots,K-1.
\end{equation}

La matrice d'attracteur est symétrisée dans l'implémentation :
\begin{equation}
W_A \leftarrow \frac{W_A + W_A^\top}{2}.
\end{equation}

\subsection{Sorties du bloc}
Le module \texttt{HippocampalScaffold.forward} retourne :
\begin{itemize}
    \item $\mathbf{h}$ : état hippocampique stabilisé (dimension \texttt{hpc\_size}) ;
    \item $\mathbf{g}$ : état concaténé des modules de grille (support structurel).
\end{itemize}

\subsection{Correspondance implémentation}
\begin{itemize}
    \item \texttt{GridCellModule.forward} : mise à jour de phase et encodage sinus/cosinus ;
    \item \texttt{HippocampalScaffold.\_init\_fixed\_projection} : initialisation gaussienne et gel de la projection ;
    \item \texttt{HippocampalScaffold.forward} : path integration $\rightarrow$ projection $\rightarrow$ attracteur.
\end{itemize}

\subsection{Détail mathématique fonction par fonction}

\paragraph{\texttt{GridCellModule.\_\_init\_\_}}
Chaque module porte:
\[
P \in \mathbb{N}^+, \qquad \boldsymbol{\phi}\in\mathbb{R}^{1\times2}.
\]
La phase $\boldsymbol{\phi}$ est stockée comme buffer persistant.

\paragraph{\texttt{GridCellModule.forward}}
Pour un batch de vitesses $U\in\mathbb{R}^{B\times2}$:
\[
\boldsymbol{\phi}\leftarrow (\boldsymbol{\phi}+U)\bmod P.
\]
Ensuite, codage trigonométrique:
\[
g(\boldsymbol{\phi})=
\left[
\sin\!\left(\frac{2\pi}{P}\boldsymbol{\phi}\right),
\cos\!\left(\frac{2\pi}{P}\boldsymbol{\phi}\right)
\right]\in\mathbb{R}^{B\times4}.
\]
Le redimensionnement dynamique de \texttt{phase} garantit la compatibilité entre batches successifs.

\paragraph{\texttt{HippocampalScaffold.\_\_init\_\_}}
Si $M$ modules:
\[
d_g = 4M,\qquad W_{g\to h}\in\mathbb{R}^{d_h\times d_g},\quad d_h=\texttt{hpc\_size}.
\]
La projection grille$\rightarrow$hippocampe est fixée après initialisation.

\paragraph{\texttt{HippocampalScaffold.\_init\_fixed\_projection}}
Initialisation gaussienne:
\[
W_{g\to h}[i,j]\sim\mathcal{N}(0,\sigma^2),
\]
puis:
\[
W_{g\to h}.\texttt{requires\_grad}\leftarrow\texttt{False}.
\]

\paragraph{\texttt{HippocampalScaffold.forward}}
Pour une entrée vitesse $U$:
\begin{align}
\mathbf{g} &= \mathrm{concat}\left(g_1(U),\dots,g_M(U)\right)\in\mathbb{R}^{B\times d_g},\\
\mathbf{h}_0 &= \mathrm{ReLU}(W_{g\to h}\mathbf{g}^\top - \theta_h)^\top.
\end{align}
Puis dynamique d'attracteur en $K$ itérations:
\[
\mathbf{h}^{(k+1)}=\tanh(\mathbf{h}^{(k)}W_A),\quad k=0,\dots,K-1.
\]
Sorties:
\[
(\mathbf{h}^{(K)},\mathbf{g}).
\]

\paragraph{\texttt{HippocampalScaffold.reset\_phases}}
Réinitialise toutes les phases:
\[
\forall m\in\{1,\dots,M\},\quad \boldsymbol{\phi}_m\leftarrow \mathbf{0}.
\]

\subsection{Interface détaillée des fonctions}

\paragraph{\texttt{GridCellModule.forward(velocity=None)}}
\begin{itemize}[leftmargin=*]
    \item \textbf{Entrée}: \texttt{velocity} de forme $(B,2)$, optionnelle.
    \item \textbf{Sortie}: état de grille $(B,4)$.
    \item \textbf{État interne}: \texttt{phase} est persistant et mis à jour à chaque appel.
\end{itemize}

\paragraph{\texttt{HippocampalScaffold.forward(velocity=None, steps=3)}}
\begin{itemize}[leftmargin=*]
    \item \textbf{Entrées}
    \begin{itemize}
        \item \texttt{velocity}: $(B,2)$;
        \item \texttt{steps}: nombre d'itérations attracteur.
    \end{itemize}
    \item \textbf{Sorties}
    \begin{itemize}
        \item \texttt{h}: état hippocampique $(B,d_h)$;
        \item \texttt{g}: état grille concaténé $(B,d_g)$.
    \end{itemize}
\end{itemize}

\paragraph{\texttt{HippocampalScaffold.reset\_phases()}}
\begin{itemize}[leftmargin=*]
    \item \textbf{Entrée}: aucune.
    \item \textbf{Effet}: remise à zéro de toutes les phases des modules de grille.
\end{itemize}

\subsection{Déroulement du calcul (pas à pas)}
\begin{enumerate}[leftmargin=*]
    \item \textbf{Mise à jour de phase}: addition vitesse + modulo période.
    \item \textbf{Codage périodique}: conversion phase $\rightarrow$ sinus/cosinus.
    \item \textbf{Concaténation multi-modules}: construction de $\mathbf{g}\in\mathbb{R}^{B\times d_g}$.
    \item \textbf{Projection fixe}: \texttt{Linear} figée grille$\rightarrow$hippocampe.
    \item \textbf{Seuil non linéaire}: \texttt{ReLU(.-hpc\_threshold)}.
    \item \textbf{Nettoyage attracteur}: itérations de \texttt{tanh(h @ W\_attractor)}.
\end{enumerate}

\subsection{Dimensions et cohérence batch}
Pour éviter les erreurs de forme lors des mini-lots variables:
\begin{itemize}[leftmargin=*]
    \item la phase interne est automatiquement redimensionnée à $B$;
    \item chaque module retourne 4 dimensions;
    \item la concaténation donne $d_g=4M$ avec $M=\lvert\texttt{grid\_periods}\rvert$.
\end{itemize}

\end{document}
